% =============================================================================
% Capítulo 2 — Fundamentação Teórica   (peso: médio)
% Fonte: plano §8. Arquitetura (estilos, ADRs, C4, atributos de qualidade); harness como middleware (R12, R21); terminologia agêntica (R8); Green AI (R2, R3, R5).
% -----------------------------------------------------------------------------
% Estrutura (decisões do autor, 17/10/2026):
%   2.1 LLMs, agentes e chamada de ferramentas   (R8, R11, Yao et al.)
%   2.2 Arquitetura: modificabilidade, acoplamento, lock-in; ADRs e C4
%   2.3 Camadas de abstração e middleware        (middleware mora aqui, não na 2.4)
%   2.4 O harness: definição e padrões           (seção já redigida, condensada)
%   2.5 Energia em inferência (Green AI)
%   2.6 Experimentação controlada
% Decisões de 17/10/2026: busca 6 executada (ADRs, C4, ISO/IEC 25010); SagaLLM (R9)
% usado na responsabilidade 3; detalhamento por responsabilidade (antiga 2.1.3)
% migrado para a Metodologia (Seção "Delimitação das responsabilidades do
% harness"); responsabilidades 4 e 5 reduzidas ao mínimo aqui.
% -----------------------------------------------------------------------------
% Seção do harness. Base: context/planning/mapeamento-harness-liu2025.md (A28).
% Decisão D5a(i), 17/10/2026: Liu et al. (R21) entra como referência PARCIAL;
% as lacunas são declaradas e completadas com outras fontes:
%   resp. 1 -> Issarny et al. (R12, middleware)
%   resp. 3 -> Chang e Geng (R9, SagaLLM), só para estado
%   resp. 4 -> R15 a R17 (padrões de resiliência; saíram da reserva)
%   resp. 5 -> documentação do OpenTelemetry (documento técnico, não revisado
%              por pares: exceção rotulada no texto, como os relatórios de mercado)
% Graus do Quadro = julgamento do autor, a aprovar (D5d).
% NÃO lidos na íntegra: Issarny et al. (2007), Stevens et al. (1974), Tyree e
% Akerman (2005), Buchgeher et al. (2023), Mavrogiorgou et al. (2025), Yao et
% al. (2023), Chang e Geng (2025) — o texto só lhes atribui o que consta do
% título, do resumo ou das notas verificadas do projeto. Lu et al. (2023a)
% aparece só como citado por Liu et al. (2025, seção 5.3), sem entrada no .bib.
% Busca 6: não se encontrou estudo controlado do efeito de ADRs ou do C4; a
% evidência é descritiva ou de mineração de repositórios. Por isso as duas
% práticas entram como meio de documentação, não como objeto de avaliação.
% Figuras sugeridas (não produzidas): (a) diagrama em camadas Aplicação ->
% harness -> provedores, com as cinco responsabilidades, derivável do C4 de
% contexto; (b) quadro "conceito -> objetivo -> uso no experimento".
% =============================================================================
\chapter{Fundamentação Teórica}
\label{cap:fundamentacao}

Este capítulo reúne os conceitos necessários para explicar as decisões do
trabalho. Parte do objeto que a aplicação consome, o modelo de linguagem usado
como agente (Seção~\ref{sec:llm-agentes}); passa pelos conceitos de arquitetura
que dão sentido ao custo de substituição do provedor
(Seções~\ref{sec:arquitetura} e~\ref{sec:abstracao}); define a camada estudada,
o \harness{} (Seção~\ref{sec:harness}); e fecha com os dois fundamentos do
experimento, a medição de energia em inferência (Seção~\ref{sec:energia}) e a
experimentação controlada (Seção~\ref{sec:experimentacao}).

% -----------------------------------------------------------------------------
\section{Modelos de linguagem, agentes e chamada de ferramentas}
\label{sec:llm-agentes}

Um modelo de linguagem de grande escala gera texto a partir de uma entrada
textual, o \textit{prompt}. Usado isoladamente, ele responde a uma pergunta por
vez; passa a ser componente de uma aplicação quando é consultado por programa,
por meio da interface de um provedor.\pendencia{REFERÊNCIA NECESSÁRIA: definição
de LLM e de modelo de fundação.}

O termo \textit{agente} tem mais de um uso na literatura. \citeonline{liu2025}
tratam como agente a IA que assume papel ativo para atingir objetivos do
usuário, capaz de perseguir objetivos e gerar planos. \citeonline{sapkota2026}
separam os \textbf{agentes de IA}, voltados à automação de tarefas específicas,
da \textbf{IA agêntica}, formada por sistemas multiagente com planejamento e
coordenação autônoma. Adota-se aqui o primeiro uso: o objeto é o agente de IA
que atua sobre uma tarefa, e não o sistema multiagente, que está fora do escopo.

O que torna um modelo capaz de agir é a chamada de ferramentas. O modelo não
executa a ferramenta: ele devolve uma solicitação de chamada, a aplicação a
executa, envia o resultado de volta, e o modelo prossegue, podendo repetir o
ciclo várias vezes até concluir.\pendencia{REFERÊNCIA NECESSÁRIA: documentação
de chamada de ferramentas dos provedores (Anthropic e OpenAI); documento
técnico, a rotular como não revisado por pares.} A intercalação de etapas de
raciocínio do modelo com ações sobre fontes externas foi proposta por
\citeonline{yao2023}. Quanto à interface entre agente e ferramenta,
\citeonline{zhu2026} comparam o protocolo \textit{Model Context Protocol} (MCP)
e o protocolo agente-agente A2A; este último trata de sistemas multiagente e,
portanto, também fica de fora.

Dessa descrição decorre a distinção que o experimento usa. Chama-se aqui
\textbf{requisição simples} a consulta ao modelo sem ferramentas, e
\textbf{requisição agêntica} a que envolve o ciclo de ferramentas, com várias
idas e vindas entre aplicação e modelo. É uma definição operacional do autor,
detalhada na Metodologia. A diferença importa por duas razões: o ciclo de
ferramentas é parte do trabalho que uma camada intermediária pode orquestrar, e
a quantidade de chamadas e o tamanho do contexto crescem a cada volta do ciclo,
o que afeta o custo de cada requisição (Seção~\ref{sec:energia}).

% -----------------------------------------------------------------------------
\section{Arquitetura de software: modificabilidade, acoplamento e dependência de fornecedor}
\label{sec:arquitetura}

\subsection{Modificabilidade, acoplamento e \textit{vendor lock-in}}
\label{subsec:arq-conceitos}

\citeonline{parnas1972} propôs que as decisões de projeto com maior
probabilidade de mudar fossem encapsuladas em módulos que as escondessem do
restante do sistema. A modificabilidade é hoje tratada como atributo de
qualidade arquitetural a ser projetado deliberadamente \cite{bass2021}.
\pendencia{Avaliar a norma ISO/IEC 25010 como base normativa para
modificabilidade e portabilidade. Não lida; a edição de 2023 reorganizou as
características (a portabilidade passou a ser chamada de flexibilidade), então
conferir a edição e os termos no texto da norma antes de citar.}
O conceito que mede o quanto uma mudança se propaga é o de acoplamento, o grau
de dependência entre módulos, que remonta a \citeonline{stevens1974}. Neste
trabalho, o código de negócio que importa o SDK do provedor e chama o modelo
diretamente está acoplado a esse provedor, e cada ponto de dependência direta é
um lugar a alterar numa troca. É o mesmo mecanismo que \citeonline{sculley2015}
descrevem como dívida técnica: o \textit{glue code} que adapta o sistema à API
de um pacote específico o prende às particularidades desse pacote.

A dependência de um único fornecedor (\textit{vendor lock-in}) é a forma
organizacional desse acoplamento: trocar de fornecedor custa tanto que a troca
deixa de ser uma opção real. \citeonline{kaur2017} revisam abordagens de
interoperabilidade e portabilidade propostas para reduzir essa dependência em
nuvens interconectadas. \citeonline{mo2023} comparam empiricamente funções
\textit{serverless} escritas sobre uma biblioteca de abstração e sobre as APIs
nativas do provedor, medindo a portabilidade e métricas de estrutura e qualidade
de código. Esse último estudo é o precedente mais próximo do desenho adotado
aqui, e a sua limitação de construto, o uso de métricas de código como
\textit{proxy} do esforço de troca, é discutida na Seção~\ref{sec:ameacas}.

O trabalho adota como definição a de \textbf{custo de substituição do
provedor} dada na Seção~\ref{sec:problema}: o esforço de alteração do código
necessário para trocar o provedor preservando o comportamento da aplicação. É a
quantidade que a modificabilidade e o acoplamento, nos termos acima, tornam
mensurável.

\subsection{Documentação arquitetural: ADRs e modelo C4}
\label{subsec:arq-documentacao}

A fase de solução do trabalho é documentada por registros de decisão
arquitetural (\textit{Architecture Decision Records}, ADRs) e pelo modelo C4.
\citeonline{tyree2005} tratam a decisão arquitetural como objeto de
documentação própria e propõem uma estrutura para registrá-la. Em estudo de
mineração de repositórios do GitHub, \citeonline{buchgeher2023} verificaram que
a adoção de ADRs ainda é baixa, embora crescente, e que a maioria dos
repositórios usa o mesmo modelo de registro. O modelo C4 descreve a arquitetura
em quatro níveis, de contexto, contêineres, componentes e código
\cite{mavrogiorgou2025}. Neste trabalho, o diagrama de contêineres é o que
mostra a posição do \harness{} entre a aplicação e os provedores, e os ADRs
registram as decisões de projeto da camada, entre elas as de recuperação de
falhas e de registro de eventos.

A literatura consultada nas buscas deste trabalho sobre as duas práticas é
descritiva ou baseada em mineração de repositórios; não se encontrou
comparação controlada do efeito delas. Por isso, ADRs e C4 são usados aqui como
meio de documentação do projeto, e não como objeto de avaliação.

% -----------------------------------------------------------------------------
\section{Camadas de abstração e \textit{middleware}}
\label{sec:abstracao}

Uma camada de abstração interpõe, entre o código que consome um serviço e o
serviço em si, uma interface própria, de modo que o consumidor passa a depender
dessa interface e não das particularidades de cada fornecedor. O
\textit{middleware} é a forma consolidada dessa ideia: \citeonline{issarny2007}
o definem como a camada de \textit{software} que oferece soluções reutilizáveis
para a heterogeneidade, a interoperabilidade, a segurança e a confiabilidade.
A heterogeneidade em jogo neste trabalho é a das convenções de cada provedor de
modelos, que definem formatos próprios de mensagens, de declaração e chamada de
ferramentas, de erros e de limites de uso (Seção~\ref{sec:contexto}).

A mesma estratégia aparece em outros pontos da literatura citada. Mitigar o
\textit{lock-in} por meio de uma camada intermediária é problema consolidado
de arquitetura fora da IA \cite{kaur2017,mo2023}. No domínio dos agentes,
\citeonline{kandogan2025} propõem uma arquitetura de referência para IA composta
com registros de agentes e de dados que abstraem provedores e APIs
proprietárias, ilustrada por um caso de uso e sem comparação controlada com uma
integração direta.

A abstração, porém, tem preço. Toda camada acrescenta indireção a cada
requisição, e o acoplamento ao provedor pode ser trocado por acoplamento à
interface da própria camada, risco que o catálogo de \citeonline{liu2025}
registra para o \textit{Tool/agent registry} (Seção~\ref{sec:delimitacao}).
Em leitura própria do catálogo, 14 dos 18 padrões listam sobrecarga
(\textit{overhead}) entre as desvantagens, e nenhum a quantifica. Esse custo da
indireção é o que a hipótese energética deste trabalho coloca à prova, e o
benefício da abstração, o que o custo de substituição mede. O conceito de
\textit{middleware} explica por que uma camada intermediária absorve a
heterogeneidade, mas não indica como projetá-la para provedores de modelos; essa
transposição é do autor. \pendencia{Issarny et al. (2007) não foi lido na
íntegra; a formulação segue a das notas do projeto. Reverificar a definição de
\textit{middleware} no texto.}

% -----------------------------------------------------------------------------
\section{O \textit{harness}: definição adotada e padrões arquiteturais}
\label{sec:harness}

\subsection{Definição adotada}
\label{subsec:harness-definicao}

O termo \harness{} é recente na literatura sobre agentes de IA e não tem
definição canônica. \citeonline{rafique2026} o empregam para a camada de
\textit{software} que envolve o modelo e se responsabiliza por montar o contexto
de cada chamada, mediar a execução de ferramentas, gerenciar o estado e observar
os eventos do ciclo de vida do agente.\pendencia{Confirmar, no texto integral de
Rafique e Bindschaedler (2026), que as quatro funções são descritas pelos
autores; o artigo trata de um sistema específico (gerência de memória virtual).}
A definição adotada neste trabalho parte
desse uso e lhe acrescenta uma posição arquitetural explícita: o \harness{} é a
\textbf{camada arquitetural intermediária entre a aplicação \textit{web} e o
provedor de IA}, que abstrai e centraliza as interações da aplicação com modelos
e agentes, de modo que a aplicação não acessa diretamente os SDKs dos
provedores. O fluxo de uma requisição passa, portanto, da aplicação ao
\harness{} e deste ao provedor. A camada não implementa regras de negócio, não
substitui a interface com o usuário e não treina modelos. Em termos da
Seção~\ref{sec:abstracao}, é um \textit{middleware} aplicado aos provedores de
modelos.

Atribuem-se à camada cinco responsabilidades: (1) abstração e padronização da
comunicação com os provedores; (2) orquestração de chamadas e execução
controlada de ferramentas; (3) gerenciamento de contexto e estado da interação;
(4) tratamento e recuperação de falhas; e (5) registro estruturado de eventos e
auditoria. Essa lista é uma definição operacional do autor, e não um resultado
da literatura. Por isso, as subseções seguintes a confrontam com publicações:
apresentam o catálogo de padrões de \citeonline{liu2025} e verificam,
responsabilidade a responsabilidade, o quanto ele a sustenta e que outra fonte
cobre o que ele não cobre.

\subsection{O catálogo de padrões de Liu et al. (2025)}
\label{subsec:harness-catalogo}

\citeonline{liu2025} apresentam 18 padrões arquiteturais para agentes baseados
em modelos de fundação, identificados a partir de uma revisão sistemática de 57
estudos e complementados, para a identificação dos usos conhecidos de cada
padrão, por literatura cinzenta e por documentação oficial de aplicações reais.
Os padrões vão da definição do objetivo e da geração de plano até a reflexão, a
cooperação entre agentes, as barreiras de proteção (\textit{guardrails}), o
registro de ferramentas e a avaliação do agente. Cada um é descrito no formato
clássico de padrão: resumo, contexto, problema, forças, solução, consequências
(benefícios e desvantagens), usos conhecidos e padrões relacionados. O escopo
declarado é o do agente capaz de perseguir objetivos e gerar planos, e o
público-alvo são os arquitetos e desenvolvedores que projetam esses agentes.

Duas características do catálogo orientam o seu uso neste trabalho. A primeira é
que ele deriva da literatura e não é uma avaliação empírica: benefícios e
desvantagens são listados de forma qualitativa, e os próprios autores observam,
na seção 5.3, que a maior parte deles são atributos de qualidade de
\textit{software} que ainda exigem quantificação. A segunda é que o foco recai
sobre o projeto interno do agente, e não sobre a integração entre uma aplicação
e o provedor do modelo, o que limita o alcance da correspondência discutida a
seguir. Os usos conhecidos que o catálogo lista, como LangChain, Semantic Kernel
e AutoGen, vêm da documentação oficial dessas ferramentas e indicam apenas que o
padrão foi aplicado nelas, e não que elas o entregam com qualidade comprovada.

\subsection{Correspondência entre as responsabilidades e a literatura}
\label{subsec:harness-correspondencia}

Classifica-se a correspondência entre cada responsabilidade e o catálogo em três
graus: \textit{parcial}, quando um padrão cobre parte da função; \textit{analogia},
quando a função é semelhante, mas o objeto é diferente; e \textit{nenhuma},
quando o catálogo não traz padrão aplicável. A classificação é um julgamento do
autor. O Quadro~\ref{qua:harness-liu} resume o resultado e indica, nas lacunas, a
fonte que as complementa.

\begin{quadro}[h!]
  \centering
  \Caption{\label{qua:harness-liu}Correspondência entre as responsabilidades do
  \harness{} e os padrões de Liu et al. (2025)}
  \UFCqua{}{
    \small
    \begin{tabularx}{\textwidth}{|>{\raggedright\arraybackslash}p{3.4cm}|>{\raggedright\arraybackslash}X|>{\raggedright\arraybackslash}p{2.1cm}|>{\raggedright\arraybackslash}X|}
      \hline
      \textbf{Responsabilidade} & \textbf{Padrões do catálogo} & \textbf{Corres\-pon\-dência} & \textbf{Complemento adotado} \\
      \hline
      1. Abstração e padronização da comunicação com os provedores
        & \textit{Agent adapter}; \textit{Prompt/response optimiser}
        & Analogia
        & Noção de \textit{middleware} \cite{issarny2007} \\
      \hline
      2. Orquestração de chamadas e execução controlada de ferramentas
        & \textit{Agent adapter}; \textit{Tool/agent registry}; \textit{One-shot} e \textit{Incremental model querying} (estes, por analogia)
        & Parcial
        & Nenhum \\
      \hline
      3. Gerenciamento de contexto e estado da interação
        & \textit{Passive goal creator}; \textit{Incremental model querying}; \textit{Retrieval augmented generation}
        & Parcial (contexto); nenhuma (estado)
        & Gestão de contexto e estado \cite{chang2025} \\
      \hline
      4. Tratamento e recuperação de falhas
        & Nenhum
        & Nenhuma
        & Padrões de resiliência \cite{mendonca2020,sedghpour2022,aderaldo2025} \\
      \hline
      5. Registro estruturado de eventos e auditoria
        & Nenhum entre os 18
        & Nenhuma
        & Documentação do OpenTelemetry \cite{opentelemetrylogs} \\
      \hline
    \end{tabularx}
  }{
    \Fonte{Elaborado pelo autor, com base em \citeonline{liu2025}}
  }
\end{quadro}

Das cinco responsabilidades, só a segunda tem correspondência parcial direta com
um padrão do catálogo; a primeira corresponde por analogia, a terceira em parte
(só quanto ao contexto), e a quarta e a quinta não têm correspondência nos 18
padrões. A razão está no escopo do catálogo: o seu objeto é o agente que
persegue objetivos e gera planos, e não a camada que o liga ao provedor do
modelo. O exame de cada correspondência, com os padrões e as fontes de
complemento, pertence ao projeto da camada e consta na
Seção~\ref{sec:delimitacao}.

Duas lacunas pedem fontes próprias. Na responsabilidade 4, o \textit{Retry} (nova
tentativa após falha transitória) e o \textit{Circuit Breaker} (interrupção
temporária das chamadas a um serviço que falha repetidamente) são padrões de
resiliência usados em microsserviços \cite{aderaldo2025}, já analisados por meio
de modelos \cite{mendonca2020}. Na responsabilidade 5, adota-se a documentação
do OpenTelemetry, projeto de padronização de telemetria, segundo a qual o
registro (\textit{log}) estruturado pode carregar o identificador do rastreio e
do segmento da requisição (\textit{TraceId} e \textit{SpanId}), que o associam à
requisição em cujo processamento foi emitido \cite{opentelemetrylogs}. Trata-se de documento técnico, não
revisado por pares, usado apenas para fixar o conceito, e não para sustentar
afirmação empírica.

\subsection{Síntese e relação com o trabalho}
\label{subsec:harness-sintese}

A definição de \harness{} adotada é operacional e de responsabilidade do autor, e
a literatura a sustenta em parte e por fontes diferentes. O catálogo de
\citeonline{liu2025} fundamenta, com as ressalvas indicadas, a responsabilidade
2 e o contexto da 3; a responsabilidade 1 apoia-se no conceito de
\textit{middleware} (Seção~\ref{sec:abstracao}); o estado, na 3, em
\citeonline{chang2025}; a 4, em padrões de resiliência de microsserviços
\cite{mendonca2020,sedghpour2022,aderaldo2025}; e a 5, em documentação técnica
de telemetria \cite{opentelemetrylogs}. Nenhuma fonte, isoladamente, sustenta a
definição inteira, e as correspondências por analogia são hipóteses de projeto, e
não resultados.

O catálogo é qualitativo, e os seus autores reconhecem que benefícios e
desvantagens dos padrões ainda exigem quantificação \cite{liu2025}. É esse espaço
que o presente trabalho ocupa, ao medir dois efeitos da camada: o custo de
substituição do provedor, associado à responsabilidade 1, e o custo energético da
camada, incluindo a orquestração de ferramentas da responsabilidade 2. As demais
responsabilidades são implementadas por serem pré-requisitos técnicos do
experimento e não são avaliadas (Seção~\ref{sec:justificativa}).

% -----------------------------------------------------------------------------
\section{Energia em inferência de modelos de linguagem}
\label{sec:energia}

A área de \textit{Green AI} propõe tratar a eficiência energética como critério
de avaliação ao lado da acurácia \cite{schwartz2020}, depois que
\citeonline{strubell2019} chamaram atenção para o custo energético e ambiental do
treinamento de modelos de linguagem. A inferência, que se repete a cada
requisição, é o foco deste trabalho. Entre as arquiteturas, as generativas de
propósito geral consomem muito mais energia que sistemas específicos: o
consumo da geração de texto, medido por \citeonline{luccioni2024}, é mais de
vinte vezes o da classificação de texto (valores na
Seção~\ref{sec:contexto}).

Duas grandezas precisam ser distinguidas. A \textbf{potência} é a taxa
instantânea de consumo, em watts; a \textbf{energia} é a potência integrada no
tempo, em joules. Uma requisição pode consumir mais energia por durar mais, com
a mesma potência, ou por exigir mais potência no mesmo intervalo. Chamam-se aqui,
respectivamente, de componente \textbf{extensional} e componente
\textbf{intensivo} do custo; a distinção operacional entre os dois, feita por
meio da potência média e da duração relatadas em separado, é do autor.
\citeonline{vellaisamy2026} decompõem a sobrecarga de orquestração executada no
processador hospedeiro, separando-a do trabalho executado no dispositivo, o que oferece um precedente
para decompor o custo de uma camada de \textit{software}, e não somente somá-lo.

Medir é necessário porque a estimativa pode errar o consumo real.
\citeonline{fernandez2025} verificaram
que estimativas baseadas em contagem de operações de ponto flutuante
subestimam o consumo real e que otimizações bem aplicadas reduzem a energia total
em até 73\%, o que mostra que o consumo depende da pilha de \textit{software}, e
não somente do modelo. Por isso o trabalho mede a energia da GPU por telemetria
do \textit{hardware}, com a biblioteca NVML, no protocolo já empregado pelo autor
\cite{barros2026}.\pendencia{REFERÊNCIA NECESSÁRIA: documentação da NVML
(documento técnico da fabricante, a rotular como não revisado por pares).}

A unidade de medida é a requisição completa, e não a chamada isolada ao modelo.
\citeonline{ifath2026} caracterizam os compromissos entre desempenho e energia em
fluxos de múltiplas requisições, entre eles o agêntico, o mesmo padrão de uso que
uma camada de \harness{} orquestra; o registro, porém, é um resumo estendido, e
não sustenta sozinho afirmação central. Quanto ao alcance, essas medições tratam
da pilha que executa a inferência. Estendê-las a uma camada de aplicação
interposta entre o código de negócio e o provedor é uma analogia, e a relevância
energética da camada permanece uma hipótese a testar (Seção~\ref{sec:contexto}).

% -----------------------------------------------------------------------------
\section{Experimentação controlada em Engenharia de Software}
\label{sec:experimentacao}

\citeonline{wohlin2024} sistematizam o experimento controlado em Engenharia de
Software: o pesquisador manipula uma ou mais variáveis independentes, mantém as
demais condições sob controle e observa o efeito sobre variáveis dependentes.
Neste trabalho, a variável independente é a forma de integração com o provedor,
com ou sem a camada de \harness{}, aplicada a duas versões funcionalmente
equivalentes de uma mesma aplicação; as variáveis dependentes são o custo de
substituição do provedor e o custo energético, medidos nos \textit{benchmarks}
descritos no Capítulo~\ref{cap:metodologia}. A equivalência funcional é o que
permite atribuir a diferença observada à camada, e não a outra causa.

Os mesmos autores classificam as ameaças à validade em quatro tipos: interna,
externa, de construto e de conclusão, taxonomia que estrutura a
Seção~\ref{sec:ameacas}. Duas decisões do desenho dependem desse referencial: o
uso de métricas de estrutura de código como medida do custo de substituição, com
precedente em \citeonline{mo2023}, e o critério de decisão final, definido antes
da coleta de dados, para que a conclusão não seja ajustada aos resultados
(Seção~\ref{subsec:objetivos-especificos}). A escolha dessas
medidas e dos procedimentos estatísticos é detalhada na Metodologia.
